\documentclass[tikz,border=4pt]{standalone}
\usepackage{ctex}
\usepackage{amsmath,amssymb}
\usetikzlibrary{arrows.meta, positioning, fit, backgrounds, shapes.geometric}
\begin{document}
\begin{tikzpicture}[
  node distance=0.5cm and 1.2cm,
  box/.style={rectangle, rounded corners=3pt, draw=blue!60, fill=blue!8,
              text centered, font=\small, minimum width=2.5cm, minimum height=0.75cm},
  formula/.style={rectangle, rounded corners=3pt, draw=orange!80, fill=orange!10,
              text centered, font=\small, minimum width=3.8cm, minimum height=0.75cm},
  arr/.style={-{Stealth[scale=1.2]}, thick},
  label/.style={font=\footnotesize, fill=white, inner sep=1pt},
]

% Input domain
\node[box] (xdomain) {输入空间 $\mathbb{R}^n$\\$\mathbf{x}$};

% Output domain
\node[box, right=3.5cm of xdomain] (ydomain) {输出空间 $\mathbb{R}^m$\\$\mathbf{y}=f(\mathbf{x})$};

% Function box
\node[formula, below=1.5cm of xdomain, xshift=2cm] (fbox) {向量函数 $f:\mathbb{R}^n\to\mathbb{R}^m$};

% Jacobian
\node[formula, below=1.2cm of fbox] (jbox) {Jacobian 矩阵 $J\in\mathbb{R}^{m\times n}$:\\$J_{ij}=\dfrac{\partial f_i}{\partial x_j}$};

% Linear approx
\node[formula, below=1.0cm of jbox] (lapprox) {线性近似：$f(\mathbf{x}+\delta)\approx f(\mathbf{x})+J\delta$};

% Arrows
\draw[arr] (xdomain) -- node[above, label]{$\mathbf{y}=f(\mathbf{x})$} (ydomain);
\draw[arr] (xdomain.south) -- ++(0,-0.6) -| (fbox.west);
\draw[arr] (ydomain.south) -- ++(0,-0.6) -| (fbox.east);
\draw[arr] (fbox) -- (jbox);
\draw[arr] (jbox) -- (lapprox);

% Dimension annotation
\node[font=\footnotesize, draw=gray, rounded corners, fill=gray!8, align=center]
  at (8.5, -1.5)
  {若 $m=1$（标量输出）：\\$J=(\nabla f)^\top$（梯度行向量）\\若 $m=n$（方阵）：\\行列式 $|\det J|$ 是换元 Jacobian};

% Special cases
\node[formula, draw=green!60!black, fill=green!8, font=\footnotesize, align=left,
      minimum width=4cm, minimum height=1.5cm]
  at (8.5, -4.5)
  {\textbf{常用场景}\\
   $f=A\mathbf{x}$：$J=A$\\
   $f=\mathbf{x}^\top A\mathbf{x}$：$\nabla f=(A+A^\top)\mathbf{x}$\\
   Softmax：$J=\mathrm{diag}(p)-pp^\top$};

\node[font=\normalsize\bfseries] at (4, 1.2) {Jacobian 矩阵作为向量函数的线性近似};

\end{tikzpicture}
\end{document}
